\chapter{Ordered collisions and two holomorphic directions}
\label{chap:ordered-collisions-residues}

Two local insertions can exchange their order along a real direction
when their holomorphic positions differ. At equal holomorphic positions,
the same exchange crosses a collision. A coefficient algebra for these
insertions must therefore retain two orders at coincidence while
identifying them away from coincidence.

This requirement gives a geometric use for complexes with multiplication.
An ordinary function records agreement between the two orders.
A complex also records their comparison on the region where exchange is
possible. The multiplication on that comparison will distinguish two
collision branches from an algebra with a nonzero element whose square
is zero.
The difference survives even when the dimensions of all resulting
cohomology groups agree.

The number of holomorphic coordinates changes this comparison.
One coordinate gives Laurent principal parts. Two coordinates give a
class on the intersection of two punctured regions, in a different
cohomological degree. We construct both coefficient systems before
introducing operations with those coefficients.

\section{Two orders and their common region}

Work over $k=\mathbb C$. Write $z$ for the difference between two
holomorphic positions on a line. Polynomial coefficients form $R=k[z]$.
On the region $z\ne0$, they form $K=k[z,z^{-1}]$.
Every element of $K$ is a finite linear combination of integer powers
of $z$. This is the localization of $R$ obtained by making $z$
invertible, as in Section~\ref{sec:cc-localization}.

Take two copies of the complex line, labelled $-$ and $+$.
Identify a point of the first copy with the point of the second copy
having the same nonzero coordinate. Keep the two origins distinct.
Polynomial functions on the two copies are pairs $(a,b)\in R^2$.
They agree on the common region precisely when $b-a=0$ in $K$.
Retaining this difference gives the cochain complex
\[
 D_1^0=R^2,\qquad D_1^1=K,\qquad d(a,b)=b-a.
\]
All unlisted degrees are zero. Its multiplication is
\begin{equation}\label{eq:ocr-ordered-cup}
 \begin{split}
 (a,b)(c,e)&=(ac,be),\\
 (a,b)v&=av,\qquad v(c,e)=ve,\qquad vw=0
 \quad(v,w\in K).
 \end{split}
\end{equation}
Thus a comparison has a left endpoint and a right endpoint.
Their actions use the corresponding components of a pair.
The unit is $(1,1)$, and $R$ acts centrally by $a\mapsto(a,a)$.

\Needspace{10\baselineskip}
\begin{proposition}\label{prop:ocr-line-algebra}
These formulas make $D_1$ an associative differential graded $R$-algebra.
Its cohomology is
\[
 H^0(D_1)=R,\qquad H^1(D_1)=K/R.
\]
Products of positive-degree cohomology classes vanish.
\end{proposition}

\begin{proof}
With three degree-zero inputs, associativity is componentwise.
With one degree-one input it follows from scalar multiplication in $K$.
For example,
$((a,b)v)(c,e)=ave=(a,b)(v(c,e))$.
With at least two degree-one inputs both parenthesizations vanish.
For two degree-zero inputs the differential identity is exactly
\[
 be-ac=(b-a)e+a(e-c).
\]
Every other Leibniz identity has zero degree-two target.
The kernel of $R^2\to K$ is the diagonal copy of $R$.
Its image is $R\subset K$, proving the cohomology statement.
\end{proof}

The quotient $K/R$ has basis $z^{-a-1}$ for integers $a\ge0$.
It retains the entire principal part of a Laurent polynomial.
The value of its residue, the coefficient of $z^{-1}$, retains only
one of these coefficients.

\section{A puncture in the complex plane of dimension two}

Let $(x,y)$ be the difference between two positions in $\mathbb C^2$.
The points are distinct if $x\ne0$ or $y\ne0$.
There are now two regions on which to compare coefficients.
Put
\[
 R=k[x,y],\quad R_x=k[x,x^{-1},y],\quad
 R_y=k[x,y,y^{-1}],\quad
 R_{xy}=k[x,x^{-1},y,y^{-1}].
\]
The subscript specifies which coordinates have been made invertible.
These are subrings of the common Laurent polynomial ring $R_{xy}$.
Agreement on the overlap is recorded by
\begin{equation}\label{eq:ocr-puncture-complex}
 B^0=R_x\oplus R_y
 \xrightarrow{\ \delta\ }B^1=R_{xy},
 \qquad \delta(u,v)=v-u.
\end{equation}
This is the two-region \v{C}ech complex: degree zero contains local
values, and degree one contains their difference on the intersection.
Give it the multiplication~\eqref{eq:ocr-ordered-cup}, using $R_x$
and $R_y$ for the endpoint rings and $R_{xy}$ for the intersection.
The same calculation proves the differential algebra identities.

\begin{proposition}\label{prop:ocr-two-dimensional-puncture}
The cohomology of $B$ is
\[
 H^0(B)=R,\qquad
 H^1(B)=J:=R_{xy}/(R_x+R_y)
       =\bigoplus_{a,b\ge0} kx^{-a-1}y^{-b-1}.
\]
The displayed direct sum consists of finite sums.
In particular, every compatible pair of Laurent polynomials on the
puncture extends to a polynomial at the origin.
\end{proposition}

\begin{proof}
A Laurent polynomial has a unique finite expansion in monomials
$x^iy^j$, with $i,j\in\mathbb Z$.
Membership in $R_x$ requires $j\ge0$ in every nonzero monomial.
Membership in $R_y$ requires $i\ge0$.
Their intersection is therefore $R$.
This computes the kernel of $\delta$.

Their sum contains exactly the Laurent polynomials whose monomials
have at least one nonnegative exponent.
Quotienting leaves precisely the monomials with both exponents
negative. Uniqueness of Laurent coefficients proves both spanning
and linear independence of the asserted basis.
\end{proof}

For example, $x^{-1}$ is defined on $x\ne0$, so it contributes zero
to $J$. The same holds for $y^{-1}$.
The fraction $x^{-1}y^{-1}$ has a nonzero class because it belongs
to neither summand, even after addition of terms from both summands.
It represents a degree-one comparison class on the puncture.
It is not a compatible degree-zero function there.

The origin can also be retained as a support condition.
Adjoin the restriction from $R$ to the two-region complex:
\begin{equation}\label{eq:ocr-support-complex}
 L^0=R\xrightarrow{\ a\mapsto(a,a)\ }
 L^1=R_x\oplus R_y
 \xrightarrow{\ (u,v)\mapsto u-v\ }L^2=R_{xy}.
\end{equation}
The same monomial calculation proves
$H^0(L)=H^1(L)=0$ and $H^2(L)=J$.
Here $J$ is the local-cohomology module at the origin, with that term
defined by the support complex~\eqref{eq:ocr-support-complex}.
Its degree is two in $L$, whereas it has degree one in $B$.

\section{Keeping the two collision branches}

Retain two copies of $\mathbb C^2$, one for each real order.
Identify their punctured planes by the coordinate identity.
The comparison on that common region is now the complex $B$,
rather than a single Laurent ring.
Its ordered gluing algebra has terms
\begin{equation}\label{eq:ocr-surface-gluing}
 \begin{split}
 D_2^0&=R^2,\\
 D_2^1&=R_x\oplus R_y,\\
 D_2^2&=R_{xy},\\
 d(a,b)&=(b-a,b-a),\\
 d(u,v)&=u-v.
 \end{split}
\end{equation}
The sign in the last line is the sign from shifting the comparison
complex by one degree. All other differentials vanish.
For a positive-degree element $w$, define
\[
 (a,b)w=aw,\qquad w(a,b)=wb.
\]
Multiply degree-zero pairs componentwise and set every product of
two positive-degree elements to zero.
The unit and the central $R$-action are diagonal.

\begin{proposition}\label{prop:ocr-surface-algebra}
The formulas above define an associative differential graded
$R$-algebra. Its only nonzero cohomology groups are
\[
 H^0(D_2)=R,\qquad H^2(D_2)=J.
\]
The induced multiplication is the square-zero extension
$R\oplus J[-2]$: the copy of $J$ lies in degree two and has
zero product with itself.
\end{proposition}

\begin{proof}
The square of the differential on $(a,b)$ is
$(b-a)-(b-a)=0$.
Associativity follows from the endpoint actions, exactly as for $D_1$.
For two degree-zero inputs the Leibniz calculation is again
$be-ac=(b-a)e+a(e-c)$ on each component of degree one.
With one positive-degree input it reduces to $R$-linearity of its
differential. With two such inputs both sides vanish.

The kernel in degree zero is the diagonal $R$.
In degree one, $u-v=0$ implies $u=v\in R_x\cap R_y=R$.
Every such pair is $d(0,u)$.
The cokernel in degree two is $R_{xy}/(R_x+R_y)=J$.
\end{proof}

The algebra has an explicit interpretation in terms of comparisons.
Let $I$ have degree-zero basis $e_-,e_+$ and degree-one basis $h$.
The nonzero products of basis elements are
\[
 e_-^2=e_-,\quad e_+^2=e_+,\quad e_-h=h=he_+,
 \qquad de_-=-h,\quad de_+=h.
\]
Its unit is $e_-+e_+$.
The two evaluation maps send the corresponding endpoint to $1$
and all other basis elements to $0$.
Both induce isomorphisms on cohomology, since
$\ker(d:I^0\to I^1)=k(e_-+e_+)$ and $d$ is surjective.

Form the differential graded tensor product $I\otimes_k B$.
Keep those elements whose two endpoint evaluations lie in
the diagonal copy of $R$ inside $B^0$.
An element has the unique form
\[
 e_-\otimes(a,a)+e_+\otimes(b,b)+h\otimes w,
 \qquad a,b\in R,\quad w\in B.
\]
For homogeneous $w$, its last term has degree $|w|+1$.
The tensor differential is
$d(h\otimes w)=-h\otimes\delta w$.
Consequently this subalgebra is precisely $D_2$.
The comparison is therefore retained together with its own
differential and both endpoint actions.

\section{Specializing to coincidence preserves multiplication}

At the origin, both $x$ and $y$ become zero.
Ordinary substitution in a complex need not preserve a
quasi-isomorphism, as the coefficient calculations in
Chapter~\ref{chap:relations-homotopies} show.
Adjoin exterior generators whose differentials are $x$ and $y$.
Let
\[
 K_R=R\otimes\Lambda(e_x,e_y),\qquad
 |e_x|=|e_y|=-1,\qquad de_x=x,\quad de_y=y.
\]
Here $\Lambda$ denotes the exterior algebra:
$e_x^2=e_y^2=0$ and $e_xe_y=-e_ye_x$.
The differential obeys the graded Leibniz rule, so
$d(e_xe_y)=xe_y-ye_x$.
This differential graded algebra is the Koszul complex of $x,y$ over $R$.

\begin{lemma}\label{lem:ocr-residue-resolution}
The augmentation $K_R\to k=R/(x,y)$ is a quasi-isomorphism.
For any quasi-isomorphism of complexes $V\to W$ of $R$-modules,
the map $V\otimes_RK_R\to W\otimes_RK_R$ is a quasi-isomorphism.
\end{lemma}

\begin{proof}
In degrees $-2,-1,0$, the complex is
\[
 R\xrightarrow{\ c\mapsto(-yc,xc)\ }R^2
 \xrightarrow{\ (a,b)\mapsto xa+yb\ }R.
\]
The first map is injective because $R$ has no zero divisors.
If $xa+yb=0$, setting $x=0$ gives $yb(0,y)=0$.
Hence $b=xc$ for a polynomial $c$, and then $a=-yc$.
This proves exactness in degree $-1$.
The cokernel in degree zero is $k$.

Filter the tensor complexes by the exterior degree in $K_R$.
The differential lowers that degree or preserves it.
There are three filtration steps, and every successive quotient
is a finite direct sum of shifts of $V$ or $W$.
The induced map on every quotient is a quasi-isomorphism.
The long exact cohomology sequences for these finite filtrations
prove the assertion.
\end{proof}

For a differential graded algebra with central $R$-action, call
$A\otimes_RK_R$ its derived fibre at the origin.
Use the tensor product multiplication with the graded signs already
fixed in Chapter~\ref{chap:relations-homotopies}.
The lemma proves that a zigzag of $R$-algebra quasi-isomorphisms
induces a zigzag between these derived fibres.
Their cohomology algebras must therefore be isomorphic.

The algebra $k\times k$ contains the elements $(1,0)$ and $(0,1)$.
Each has square equal to itself, their product is zero, and their sum
is the unit. An element with $e^2=e$ is called an idempotent.
Two idempotents are orthogonal when their products in both orders
are zero. These two idempotents distinguish the collision branches.
An element is nilpotent when some positive power of it is zero.

\begin{theorem}\label{thm:ocr-relative-nonformality}
The ordered gluing algebra $D_2$ is not formal over $R$.
Its derived fibre has degree-zero algebra $k\times k$.
The derived fibre of its cohomology algebra has degree-zero algebra
$k[\epsilon]/(\epsilon^2)$.
Both fibres have zero cohomology in every other degree.
\end{theorem}

\begin{proof}
Each localization of $R$ is flat.
Indeed, localization of a module sends a class to zero precisely
when some denominator annihilates it.
This criterion preserves both injections and exactness at the middle
of an exact sequence, and fractions preserve surjections.
Tensoring by the localized ring is this localization of modules.
Thus every term of $D_2$ is flat.

For a bounded complex of flat modules, tensoring the augmentation
$K_R\to k$ is a quasi-isomorphism.
Filter by the finitely many degrees of that bounded complex.
On each quotient, flatness preserves the exact augmented resolution
of the preceding lemma. The finite filtration proves the claim.
It follows that $D_2\otimes_RK_R\to D_2\otimes_Rk$ is a
quasi-isomorphism of algebras.
Each localization term becomes zero: an invertible coordinate
cannot act by zero on a nonzero module.
Only $k^2$ in degree zero remains, with componentwise multiplication.

Construct a flat model for the cohomology algebra by setting
\[
 E=R\oplus L,
 \qquad (a,\ell)(b,m)=(ab,am+\ell b),
 \qquad d(a,\ell)=(0,d_L\ell).
\]
The entire summand $L$ has square-zero multiplication.
The map $E\to R\oplus J[-2]$ is identity on $R$, the quotient
map $R_{xy}\to J$ in degree two, and zero on the other terms of $L$.
It is multiplicative and is a quasi-isomorphism by the calculation
of $H^*(L)$.
The terms of $E$ are again flat and bounded.
Therefore its derived fibre is computed by $E\otimes_Rk$.
Every localization vanishes, leaving the first $k$ and the degree-zero
copy of $k$ from $L^0$.
The latter is square zero. This gives $k[\epsilon]/(\epsilon^2)$.

The element $\epsilon$ is nonzero and nilpotent.
The algebra $k\times k$ has no nonzero nilpotent element, since
$(a,b)^n=0$ forces $a=b=0$ in a field.
The two fibre algebras cannot be isomorphic.
The preceding lemma therefore excludes every zigzag of
$R$-algebra quasi-isomorphisms from $D_2$ to its cohomology.
\end{proof}

The same argument in one coordinate uses
$k[z]\otimes\Lambda(e_z)$ with $de_z=z$.
Replace $L$ by $[R\to K]$ in degrees zero and one.
It proves the corresponding nonformality of $D_1$ over $k[z]$.

The collision branches are visible before any choice of generators
for a state algebra. Specialization retains their two orthogonal
idempotents. Replacing the coefficient complex by its cohomology
changes that multiplication into a square-zero extension.
Thus higher multiplication data in the coefficient algebra can be
forced by the geometry of exchanging insertions.

\begin{figure}[H]
\centering
\[
\begin{array}{c@{\qquad\qquad}c}
 D_2 & H^*(D_2)=R\oplus J[-2]\\[4pt]
 \Big\downarrow\scriptstyle{-\otimes_RK_R}
 &\Big\downarrow\scriptstyle{-\otimes_RK_R}\\[4pt]
 k\times k & k[\epsilon]/(\epsilon^2)\\[4pt]
 (1,0)^2=(1,0),\quad(0,1)^2=(0,1)
 & \epsilon\ne0,\quad\epsilon^2=0
\end{array}
\]
\caption{At coincidence, two independent order branches and a square-zero
class have the same vector-space dimension. Their products distinguish
the two derived fibres.}
\end{figure}

\section{Residues pair singular coefficients with formal jets}

To measure a coefficient near a fixed point, retain all its Taylor
coefficients. The completed ring is
\[
 \widehat R=k[[x,y]],\qquad \mathfrak m=(x,y).
\]
The $\mathfrak m$-adic topology has $\mathfrak m^n$ as neighborhoods
of zero. A $k$-linear functional $\lambda:\widehat R\to k$, with
$k$ discrete, is continuous exactly when it kills some $\mathfrak m^n$.
This is the continuity criterion of Section~\ref{sec:cc-truncations}.
Equivalently, $\lambda$ depends on finitely many Taylor coefficients.

Repeat the localization quotient over $\widehat R$:
\[
 \widehat J=
 \widehat R[x^{-1},y^{-1}]/
 \bigl(\widehat R[x^{-1}]+\widehat R[y^{-1}]\bigr).
\]
Every numerator divided by $x^py^q$ has a unique doubly negative
part, involving only the terms with powers of $x$ below $p$ and
powers of $y$ below $q$ in the numerator.
That is a finite rectangle of coefficients.
The remaining terms lie in the two submodules being divided out.
Thus $\widehat J$ has the same $k$-basis as $J$.

Attach the oriented symbol $dx\wedge dy$, with
$dy\wedge dx=-dx\wedge dy$, and put
\[
 \rho_{ab}=x^{-a-1}y^{-b-1}\,dx\wedge dy,
 \qquad a,b\ge0.
\]
It records the chosen ordering of the two coordinates in the residue.
Define $\operatorname{Res}$ to extract the coefficient of
$x^{-1}y^{-1}\,dx\wedge dy$.
Multiplication by a formal series is well-defined on every $\rho_{ab}$:
only finitely many of its Taylor coefficients survive the quotient.

\begin{proposition}\label{prop:ocr-residue-dual}
The pairing
\[
 \langle f,\rho\rangle=\operatorname{Res}(f\rho)
\]
identifies $\widehat J\,dx\wedge dy$ with the continuous $k$-linear
dual of $\widehat R$. Explicitly,
\[
 \langle x^iy^j,\rho_{ab}\rangle=\delta_{ia}\delta_{jb}.
\]
Here $\delta_{ia}$ is $1$ when $i=a$ and $0$ otherwise.
\end{proposition}

\begin{proof}
Multiplication changes the powers of $\rho_{ab}$ to
$i-a-1$ and $j-b-1$.
Their residue is nonzero precisely when $i=a$ and $j=b$.
The functional associated to $\rho_{ab}$ therefore kills
$\mathfrak m^{a+b+1}$ and is continuous.
Every finite linear combination remains continuous.

Conversely, a continuous functional factors through
$\widehat R/\mathfrak m^n$ for some $n$.
This quotient has finite basis $x^iy^j$ with $i+j<n$.
The displayed pairing represents its dual basis by the corresponding
$\rho_{ij}$. This proves surjectivity.
The same pairing proves injectivity coefficient by coefficient.
\end{proof}

Multiplication by $x$ lowers the first index:
\[
 x\rho_{ab}=\rho_{a-1,b}\quad(a\ge1),\qquad x\rho_{0b}=0.
\]
There is an identical formula for $y$ and the second index.
Formal differentiation gives
\[
 \partial_x\rho_{ab}=-(a+1)\rho_{a+1,b},\qquad
 \partial_y\rho_{ab}=-(b+1)\rho_{a,b+1}.
\]
These operations obey $[\partial_x,x]=1$ and
$[\partial_y,y]=1$, with every cross commutator zero.
They follow either by differentiating Laurent monomials or by direct
substitution in the displayed formulas, including $a=0$ and $b=0$.
In characteristic zero,
\[
 \rho_{ab}=\frac{(-1)^{a+b}}{a!b!}
 \partial_x^a\partial_y^b\rho_{00}.
\]
Thus a single supported coefficient generates all derivative
coefficients, with their normalization fixed.

The scalar residue by itself is not $\widehat R$-linear to
$k=\widehat R/\mathfrak m$.
Indeed, $\operatorname{Res}(x\rho_{10})=1$, while $x$ acts as zero
on $k$. The whole principal-part module is required when an
operation must retain the action of coordinate functions.

\section{Derived restriction to a coordinate line}

Consider the coordinate line $y=0$.
Setting $y=0$ directly in a fraction containing $y^{-1}$ is undefined.
Even the quotient operation on $J$ removes all classes, since
multiplication by $y$ is surjective:
$y\rho_{a,b+1}=\rho_{ab}$.
The ordinary quotient $J/yJ$ is zero.

Derived restriction to $y=0$ tensors over $k[y]$ with the Koszul
complex $k[y]\otimes\Lambda(e_y)$, $de_y=y$.
On the coefficient module $J$, this gives
\[
 J e_y\xrightarrow{\ y\ }J,
\]
in degrees $-1$ and zero.
Its degree-zero cohomology vanishes.
Its degree-minus-one cohomology is
\[
 \ker(y:J\to J)=\bigoplus_{a\ge0} kx^{-a-1}y^{-1}.
\]
The residue in $y$, normalized by
$x^{-a-1}y^{-1}\,dx\wedge dy\mapsto x^{-a-1}\,dx$,
identifies this kernel with one-variable principal parts.
The degree-minus-one Koszul generator places this kernel one degree
below the ordinary quotient.

For the supported coefficient $J[-2]$, the same calculation places
the resulting one-variable principal parts in degree one.
This matches their degree in the one-variable support complex.
A restriction of an operator algebra must also carry its
multiplication and brackets through this comparison.
The coefficient calculation gives the map and its degree, which
those operations must respect.

\section{Commuting matrices and their Koszul complex}

Fix an integer $N\ge1$ and introduce two $N\times N$ matrices
$X=(X^i_j)$ and $Y=(Y^i_j)$ of commuting polynomial variables.
Their ordinary matrix products need not commute.
The equation for a commuting pair is
\[
 (XY-YX)^i_j=\sum_{a=1}^N
 (X^i_aY^a_j-Y^i_aX^a_j)=0.
\]

For matrices $X,Y$ with entries in $k$, this equation has an elementary
module interpretation.
A $k[x,y]$-module structure on $k^N$ extending its scalar action
assigns commuting linear maps to multiplication by $x$ and $y$.
Conversely, commuting matrices define the action
\[
 \left(\sum_{a,b\ge0}c_{ab}x^ay^b\right)v
   =\sum_{a,b\ge0}c_{ab}X^aY^bv,\qquad v\in k^N,
\]
where $c_{ab}\in k$ and each sum is finite.
For nonnegative integers $a,b,c,d$, commutation gives
$X^aY^bX^cY^d=X^{a+c}Y^{b+d}$.
This proves the multiplicative module law. The identity matrix supplies the unit law,
and addition and scalar multiplication follow term by term.
Thus the matrix-entry equations describe these module structures in
a fixed basis.

Adjoin a degree-minus-one exterior generator $\psi^i_j$ for each entry.
Define the commutative graded algebra
\[
 \mathcal R_N=k[X^i_j,Y^i_j\mid 1\le i,j\le N]
 \otimes\Lambda(\psi^i_j\mid1\le i,j\le N).
\]
The differential is
\[
 d\psi^i_j=(XY-YX)^i_j,
 \quad dX^i_j=dY^i_j=0.
\]
The exterior generators anticommute.
Extend $d$ by the graded Leibniz rule.
Its square vanishes on every generator and hence on every product.
This is the Koszul complex of the matrix entries of $XY-YX$.
It defines the derived commuting-pair equation, including its
negative-degree cohomology.

The trace of a matrix is the sum of its diagonal entries.
Commutativity of the coefficient ring gives
\[
 \operatorname{Tr}(XY-YX)
 =\sum_{i,a}(X^i_aY^a_i-Y^i_aX^a_i)=0,
\]
by interchanging $i$ and $a$ in the second sum.
Consequently $\operatorname{Tr}\psi=\sum_i\psi^i_i$ is a cycle.

\begin{proposition}\label{prop:ocr-scalar-brane-class}
For every $N\ge1$, the class $[\operatorname{Tr}\psi]$ is nonzero
in $H^{-1}(\mathcal R_N)$.
At $N=1$ the entire differential vanishes and
\[
 H^*(\mathcal R_1)=k[X,Y]\otimes\Lambda(\psi).
\]
\end{proposition}

\begin{proof}
Set every entry of $X$ and $Y$ to zero.
This defines a map of differential graded algebras
$\mathcal R_N\to\Lambda(\psi^i_j)$ with zero differential in the
target. The image of $\operatorname{Tr}\psi$ is the nonzero sum
of distinct exterior generators $\psi^i_i$.
A boundary would map to zero in a complex with zero differential.
The asserted class therefore cannot be a boundary.
For $N=1$ the commutator is identically zero, proving the last claim.
\end{proof}

Choosing the ordinary coordinate algebra $k[X,Y]$ at rank one
amounts to additionally setting $\psi=0$.
The calculation identifies precisely which class that choice removes.
The generators are primitives of the commuting-matrix equations, as
$\tau$ is a primitive of the scalar relation in
Chapter~\ref{chap:relations-homotopies}.
Here the generators anticommute and square to zero, as required by
the displayed exterior algebra.
